% SC26 AD template, jennfshr/sc26-repro b5195e67d9ad0b5d07e8b6840558c7251c73b3c0.
% Included after the paper's references using the supplied IEEEtran/sc26repro style.
% No AE appendix or badge claim is submitted.
%\usetag{explanation}
%\usetag{example}
\appendixAD

\section{Overview of Contributions and Artifacts}

\subsection{Paper's Main Contributions}

\begin{description}
\item[$C_1$] A numerically checked comparison of a compact Vulkan/Metal
training and inference stack with compiled PyTorch on six GPUs from four
vendors, including failures and models up to 1.7B parameters.
\item[$C_2$] Measurements of compile-time graph and kernel search: what it
explores, what it costs, and what it gains over no search and a two-second
tuner at the same revision.
\item[$C_3$] Deployment evidence: Vulkan qualification on three integrated
GPUs, and an Android XR application that shares its Vulkan queue with
rendering.
\end{description}

\subsection{Computational Artifacts}

\begin{description}
\item[$A_1$] The public
\href{https://github.com/kvark/meganeura/tree/0dbfcc0029bf98b33a03fc792e1f4e90ead17f23}{Meganeura}
and \href{https://github.com/kvark/inferena/tree/7b8fcb72e55410d8e33bbe948f89be50c691fcf7}{Inferena}
sources, both tagged \code{paper-p3hpc-2026-final}, and the paper's
supplement: all campaign and search-ablation records, the ablation runner
patch, and CPU-only analyzers. No DOI has been assigned.
\item[$A_2$] The public
\href{https://github.com/kvark/dinovision/tree/dc35cdf1c7c910cdd93c5b5362846842ae469a21}{DinoVision source}
and \href{https://huggingface.co/mad-bot/dinovision/tree/2c3f9017fe74c41482b165890c14737a2ccd4b6a}{evidence archive}.
No DOI has been assigned.
\end{description}

\begin{center}
\begin{tabular}{lll}
\toprule
Artifact & Contributions & Related paper elements\\
\midrule
$A_1$ & $C_1$, $C_2$, $C_3$ & Tables~\ref{tab:devices}--\ref{tab:preparation}, Fig.~\ref{fig:ratios}\\
$A_2$ & $C_3$ & Section~\ref{sec:edge}, Fig.~\ref{fig:dinovision}\\
\bottomrule
\end{tabular}
\end{center}

\section{Artifact Identification}

\newartifact

\artrel
The sources implement the compiler, validation tools and paired collection.
The supplement preserves the campaigns used in every framework timing table and
in Fig.~\ref{fig:ratios}, and the native-only search ablation. Separate
driver bring-up notes document availability limits; they supply no timing.

\artexp
Offline analysis should reproduce the generated tables and figure exactly.
It checks 198 complete GPU pairs, six inference-only pairs with separate
eager diagnostics, 30 CPU-oracle qualifications, and 60 ablation processes
against the main campaigns' PyTorch outputs. Two hard-failed campaigns are retained,
not replaced by successful retries. Fresh runs reproduce the procedure, not
the original timing noise or the exact driver-dependent failures.

\arttime
Offline setup and analysis each take about 1--5 minutes on an ordinary
workstation; no GPU is needed. For fresh collection, allow roughly
10--30 minutes for setup plus downloads, and 120--240 minutes per GPU
campaign. The search ablation took about 20 minutes for two GPUs. These are
planning estimates, not measured end-to-end costs. Large models and
CPU-oracle qualification can take longer. Repeating all platforms requires
separate machines; it is not a single short run.

\artin

\artinpart{Hardware}
Offline replay needs only a CPU and enough disk for the supplied records.
Fresh collection needs the relevant Vulkan or Metal GPU plus its
PyTorch backend. The primary systems are RTX~5070, RTX~3050, H100,
RX~7900~XT, Arc~B570 and Apple~M3. The qualification-only GPUs are
Radeon~780M, Ryzen~9600X's integrated Radeon and Intel RPL-U.
H100 has 80\,GB device memory for the 1.7B extension.
Drivers and replay differences are in Table~\ref{tab:devices}
and the campaign manifests.

\artinpart{Software}
Analysis uses standard-library Python 3.11 or later.
Collection uses Python 3.13.13 and PyTorch 2.13.0, source
\code{cf30153c4c131c8164ee7798e5022d810682e2cb}, from the
\href{https://download.pytorch.org/whl/}{PyTorch wheel index}.
Inferena pins Meganeura \code{0dbfcc00}, Blade \code{fbb4f28},
the Rust dependency lockfile, and platform-specific Python requirements.
The measured Inferena revision is \code{7b8fcb72}; the search ablation adds a
runner patch that only exposes two existing library options.
The pinned \href{https://github.com/kvark/inferena/blob/7b8fcb72e55410d8e33bbe948f89be50c691fcf7/EXPERIMENT.md}{collection instructions}
cover Linux, Windows and macOS, including the Windows Rust build toolchain.
Manifests retain package versions and PyTorch build configuration.
Rust compiler versions were not retained in every campaign receipt;
this limits exact reconstruction of native builds.

\artinpart{Datasets / Inputs}
Inputs are synthetic and deterministic. Non-language-model weights use
\code{name-index-uniform-v1}. The setup downloads SmolLM2 checkpoints
from \href{https://huggingface.co/HuggingFaceTB}{HuggingFaceTB} at pinned
model revisions; the manifest records their hashes.
The supplement contains all retained JSON values and 244 text logs from
fourteen campaign archives and two standalone crash logs, compressed
as \path{records.jsonl.xz}, and \path{search-ablation.tgz} with the 60
ablation records, their logs, driver script, runner patch and kernel-level
tile study. It excludes weights, executables and caches.
Archive hashes establish provenance; the supplement's own manifest
protects the packaged files.

\artinpart{Installation and Deployment}
For offline analysis, unpack the supplement and use its README.
For fresh collection, check out the Inferena tag above, install Rust
and \href{https://docs.astral.sh/uv/}{uv}, and run the platform setup
script with \code{cu130}, \code{rocm7.2}, \code{xpu}, \code{mps} or
\code{cpu}. It creates an isolated environment with managed Python.
Use the Windows setup script on Windows.
The collector builds the native runner and prepares pinned models.
Driver installation remains a platform prerequisite, not a Python
environment operation.

\artcomp
Setup precedes device selection and collection; collection produces
manifests and per-process records; CPU-only analysis validates those
records and generates the tables.
Run \code{python scripts/p3hpc.py} in the managed environment for the
five-model matrix. It uses three fresh process pairs per condition,
both arithmetic contracts, at least five warmups and two seconds per
phase, then twenty samples. Native search has a shared soft 60-second
deadline per session; default PyTorch compilation has a supervised
120-second deadline. CUDA/HIP/XPU use checked replay; MPS compiles
without a public equivalent.
Use \code{--models SmolLM2-360M} for the desktop extension and add
\code{SmolLM2-1.7B} on H100. CPU-oracle qualification uses
\code{--backend cpu --qualify-only --eager}, never a CPU/GPU speed ratio.
\code{--gpu} selects the native GPU on multi-GPU hosts.
Copy or rename \code{latest.tgz} before another campaign replaces it.
For the search ablation, apply \path{ablation-runner.patch} to Inferena
\code{7b8fcb72} and run \path{ablation.py} from the archive on the GPU host.

\artout
Run \code{python cohort.py records.jsonl.xz --check tables --output regenerated},
then \code{python search\_ablation.py records.jsonl.xz search-ablation.tgz --check tables --output regenerated}.
The first analyzer checks source identity, raw/joined agreement, numerical
gates, preparation and replay policies, and phase-level replication; the
second checks every ablation record's identity and policy and validates its
outputs against the main campaigns' PyTorch outputs. Together they regenerate nine
LaTeX fragments, including the ratio figure, plus per-condition CSV medians
and ranges, failure records and search summaries. Failed PyTorch phases
remain absent; independently qualified native timings remain available.
Eager diagnostics never replace compiled timing. The supplement README
explains the M3 timing sensitivity and the six-platform portability scores.
The artifact has not undergone independent artifact evaluation.

\newartifact

\artrel
DinoVision exercises native checkpoint transfer and shared rendering/ML
resources on a physical headset, separately from the framework comparison.

\artexp
The retained outputs should reproduce the stated host/device error checks,
and the sweep records should reproduce the reported queue-sharing tradeoffs.
It is not a PyTorch comparison or an on-device training benchmark.

\arttime
Allow about 5--15 minutes to inspect the retained records. A new physical
run requires Android deployment and headset setup; allow 60--120 minutes
if the Android toolchain is already installed. These are estimates.

\artin

\artinpart{Hardware}
Offline inspection needs only a CPU. Physical reproduction requires a
Quest~3S, a development host and a GPU for decoder training.

\artinpart{Software}
The pinned DinoVision source contains the Rust dependency lockfile and
deployment instructions. Android and OpenXR tools are needed for physical
deployment. The archive records the measured application and device setup;
no Android framework comparison is claimed.

\artinpart{Datasets / Inputs}
The public evidence archive contains the fixed frame, parameters, output
checks, raw sweep samples, manifests and full-resolution media.

\artinpart{Installation and Deployment}
Follow the pinned source and archive instructions to build and install the
application. Offline analysis does not require the headset or a new build.

\artcomp
Train and save the decoder on the host, transfer its checkpoint, then run
the joined encoder/decoder on the headset's rendering context.
Check the fixed-frame host/device outputs before the live queue-sharing
sweep. Retained evidence allows these two checks to be inspected separately.

\artout
The output check measures numerical agreement, not image-quality superiority.
The sweep measures application submissions, eye updates and worker latency.
Its records support Section~\ref{sec:edge}; none enters the framework timing
tables or portability aggregate.
